RNA underpins many of the most transformative advances in \ modern biotechnology, including RNA interference, CRISPR systems, metabolite-sensing riboswitches, and recently discovered Bridge RNA-guided recombination systems. 
Despite their importance, discovering new functional RNAs remains challenging. 
Unlike proteins, RNA sequences often evolve rapidly and are conserved primarily through structure rather than sequence, limiting the effectiveness of conventional homology-search methods. 
In contrast, protein homology detection is a more tractable problem with robust and fast implementations.

Here, we present a protein-first RNA discovery framework that leverages conserved bacterial protein domains to identify candidate structured RNAs within neighbouring non-coding regions. 
We use comparative genomics and identify covarying sites, which mutate together across evolution to preserve base pairing, as evidence of real structure. 
As an initial application, we focus on light-associated protein systems to search for novel light-responsive RNA regulators or RNA-protein complexes, in effect looking for potential RNA light switches. 
Two weak BLUF-associated candidates were identified, and the FMN riboswitch positive control was rediscovered, while the negative controls yeilded no strong hits.
We further introduce an approach for measuring relatedness between RNA elements, addressing a gap in this area. 
Together, this framework offers a targeted and potentially scalable route to RNA discovery, allowing biological pathways of interest to be screened for previously uncharacterized functional RNA elements.